\
\begin{filecontents*}{refs.bib}
@inproceedings{Austin2024PreTeXt,
  author    = {Austin, David and Beezer, Robert and Cantino, Michael and Kolesnikov, Alexei and Maneki, Al and Sorge, Volker},
  title     = {PreTeXt as Authoring Format for Accessible Alternative Media},
  booktitle = {Computers Helping People with Special Needs (ICCHP 2024), Part I},
  series    = {Lecture Notes in Computer Science},
  volume    = {14750},
  pages     = {184--191},
  publisher = {Springer},
  year      = {2024},
  doi       = {10.1007/978-3-031-62846-7_22},
  url       = {https://doi.org/10.1007/978-3-031-62846-7_22}
}

@article{Doore2023MDS,
  author  = {Doore, Stacy A. and Dimmel, Justin and Kaplan, Toni M. and Guenther, Benjamin A. and Giudice, Nicholas A.},
  title   = {Multimodality as Universality: Designing Inclusive Accessibility to Graphical Information},
  journal = {Frontiers in Education},
  volume  = {8},
  pages   = {1071759},
  year    = {2023},
  doi     = {10.3389/feduc.2023.1071759},
  url     = {https://doi.org/10.3389/feduc.2023.1071759}
}

@inproceedings{Sharif2021DataVis,
  author    = {Sharif, Ather and Chintalapati, Sanjana S. and Wobbrock, Jacob O. and Reinecke, Katharina},
  title     = {Understanding Screen-Reader Users' Experiences with Online Data Visualizations},
  booktitle = {Proceedings of the 23rd International ACM SIGACCESS Conference on Computers and Accessibility (ASSETS '21)},
  pages     = {14:1--14:16},
  publisher = {ACM},
  year      = {2021},
  doi       = {10.1145/3441852.3476549},
  url       = {https://doi.org/10.1145/3441852.3476549}
}

@inproceedings{Carruthers2023PlantUML,
  author    = {Carruthers, Sarah and Thomas, Amber and Kaufman-Willis, Liam and Wang, Aaron},
  title     = {Growing an Accessible and Inclusive Systems Design Course with PlantUML},
  booktitle = {Proceedings of the 54th ACM Technical Symposium on Computer Science Education (SIGCSE '23)},
  pages     = {1007--1013},
  publisher = {ACM},
  year      = {2023},
  doi       = {10.1145/3545945.3569786},
  url       = {https://doi.org/10.1145/3545945.3569786}
}

@inproceedings{Petrausch2016UML,
  author    = {Petrausch, Vanessa and Seifermann, Stephan and M{\"u}ller, Karin},
  title     = {Guidelines for Accessible Textual UML Modeling Notations},
  booktitle = {Universal Access in HCI: Methods, Techniques, and Best Practices (UAHCI/HCII 2016), Part I},
  series    = {Lecture Notes in Computer Science},
  volume    = {9737},
  pages     = {67--74},
  publisher = {Springer},
  year      = {2016},
  doi       = {10.1007/978-3-319-41267-2_10},
  url       = {https://doi.org/10.1007/978-3-319-41267-2_10}
}

@inproceedings{Serin2022Diagrams,
  author    = {Serin, Fr{\'e}d{\'e}ric and Romeo, Katerine},
  title     = {Drawing and Understanding Diagrams: An Accessible Approach Dedicated to Blind People},
  booktitle = {Universal Access in HCI: User and Context Diversity (UAHCI/HCII 2022), Part II},
  series    = {Lecture Notes in Computer Science},
  volume    = {13309},
  pages     = {94--109},
  publisher = {Springer},
  year      = {2022},
  doi       = {10.1007/978-3-031-05039-8_7},
  url       = {https://doi.org/10.1007/978-3-031-05039-8_7}
}

@article{Wabinski2022TactileMaps,
  author  = {Wabi{\'n}ski, Jakub and Mo{\'s}cicka, Albina and Touya, Guillaume},
  title   = {Guidelines for Standardizing the Design of Tactile Maps: A Review of Research and Best Practice},
  journal = {The Cartographic Journal},
  volume  = {59},
  number  = {3},
  pages   = {239--258},
  year    = {2022},
  doi     = {10.1080/00087041.2022.2097760},
  url     = {https://doi.org/10.1080/00087041.2022.2097760}
}

@article{Barros2023Tactile3D,
  author  = {Barros, Gutenberg and Correia, Walter and Teixeira, Jo{\~a}o and de Melo, Leonardo and de Melo, Helio B. P.},
  title   = {Towards the Effectiveness of 3D Printing on Tactile Content Creation for Visually Impaired Users},
  journal = {Polymers},
  volume  = {15},
  number  = {9},
  pages   = {2180},
  year    = {2023},
  doi     = {10.3390/polym15092180},
  url     = {https://doi.org/10.3390/polym15092180}
}

@inproceedings{Clepper2025Tactile,
  author    = {Clepper, Gina and McDonnell, Emma J. and Findlater, Leah and Peek, Nadya},
  title     = {``What Would I Want to Make? Probably Everything'': Practices and Speculations of Blind and Low Vision Tactile Graphics Creators},
  booktitle = {Proceedings of the CHI Conference on Human Factors in Computing Systems (CHI '25)},
  pages     = {1159:1--1159:16},
  publisher = {ACM},
  year      = {2025},
  doi       = {10.1145/3706598.3714173},
  url       = {https://doi.org/10.1145/3706598.3714173}
}

@inproceedings{Holloway2022InfoSonics,
  author    = {Holloway, Lauren M. and Goncu, Cem and Ilsar, Adrien and Butler, Mark and Marriott, Kim},
  title     = {InfoSonics: Accessible Infographics for People Who Are Blind Using Sonification and Voice},
  booktitle = {Proceedings of the CHI Conference on Human Factors in Computing Systems (CHI '22)},
  pages     = {480:1--480:13},
  publisher = {ACM},
  year      = {2022},
  doi       = {10.1145/3491102.3517465},
  url       = {https://doi.org/10.1145/3491102.3517465}
}

@inproceedings{Alam2023SeeChart,
  author    = {Alam, Md. Zakirul I. and Islam, Sk. Mainul and Hoque, Enamul},
  title     = {SeeChart: Enabling Accessible Visualizations through Interactive Natural Language Interface for People with Visual Impairments},
  booktitle = {Proceedings of the 28th International Conference on Intelligent User Interfaces (IUI '23)},
  pages     = {46--64},
  publisher = {ACM},
  year      = {2023},
  doi       = {10.1145/3581641.3584099},
  url       = {https://doi.org/10.1145/3581641.3584099}
}

@article{Seo2019RMarkdown,
  author  = {Seo, JooYoung and McCurry, Sean},
  title   = {LaTeX is NOT Easy: Creating Accessible Scientific Documents with R Markdown},
  journal = {Journal on Technology and Persons with Disabilities},
  volume  = {7},
  pages   = {157--171},
  year    = {2019}
}

@inproceedings{Armano2018LaTeXPDF,
  author    = {Armano, Tiziana and Capietto, Anna and Coriasco, Sandro and Murru, Nadir and Ruighi, Alessandro and Taranto, Enrico},
  title     = {An Automatized Method Based on LaTeX for the Realization of Accessible PDF Documents Containing Formulae},
  booktitle = {Computers Helping People with Special Needs (ICCHP 2018), Part II},
  series    = {Lecture Notes in Computer Science},
  volume    = {10896},
  pages     = {583--589},
  publisher = {Springer},
  year      = {2018},
  doi       = {10.1007/978-3-319-94277-3_91},
  url       = {https://doi.org/10.1007/978-3-319-94277-3_91}
}

@article{Leotta2023AutoCaptioning,
  author  = {Leotta, Francesco and others},
  title   = {Evaluating the Effectiveness of Automatic Image Captioning for Web Accessibility},
  journal = {Universal Access in the Information Society},
  year    = {2023},
  note    = {Bitte DOI/Seiten bei Verf{\"u}gbarkeit nachtragen.}
}

@article{Edwards2023AltTextProcess,
  author  = {Edwards, W. and others},
  title   = {How the Alt Text Gets Made: Roles and Processes of Alternative Text Creation},
  journal = {ACM Transactions on Computer-Human Interaction},
  year    = {2023},
  note    = {Bitte DOI/Seiten bei Verf{\"u}gbarkeit nachtragen.}
}

@article{Jung2021AltText,
  author  = {Jung, Crescentia and Mehta, Shubham and Kulkarni, Atharva and Zhao, Yuhang and Kim, Yea-Seul},
  title   = {Communicating Visualizations without Visuals: Investigation of Visualization Alternative Text for People with Visual Impairments},
  journal = {arXiv e-prints},
  eprint  = {2108.00141},
  year    = {2021},
  note    = {Preprint}
}
\end{filecontents*}

\documentclass[a4paper,11pt]{article}
\usepackage[utf8]{inputenc}
\usepackage[T1]{fontenc}
\usepackage{lmodern}
\usepackage[ngerman]{babel}
\usepackage{csquotes}
\usepackage{hyperref}
\usepackage[left=2.5cm,right=2.5cm,top=2.5cm,bottom=3cm]{geometry}

\title{Priorisierte Literatur-Recherche (EuroPLoP 2025)\\ \large Quellenempfehlungen pro Abschnitt mit Kurzbegründung}
\author{}
\date{\today}

\begin{document}
\maketitle

\section*{Hinweis}
Dieses Dokument ist eigenständig kompilierbar und enthält alle Bib\TeX{}-Einträge im Kopf (via \texttt{filecontents}). Die Auswahl priorisiert belastbare, aktuelle Arbeiten (2016--2025) zu Barrierefreiheit in \emph{CS/EE}-Lehre für BVI/BLV. 

\section{Paper allgemein (übergreifend)}
\begin{enumerate}
  \item \textbf{Austin et al. (2024)} \cite{Austin2024PreTeXt}.\\
  \emph{Beschreibung:} Stellt \emph{PreTeXt} als quellformat-zentrierten Workflow vor, der aus einer textuellen Quelle barrierefreie Ausgaben (HTML, PDF, EPUB, Braille) erzeugt; demonstriert STEM-Beispiele.\\
  \emph{Begründung:} Stützt euer \emph{text-first}/transformierbares Vorgehen über das gesamte Paper hinweg (Core-Content $\rightarrow$ Multi-Output) und liefert zitierfähige Evidenz für ein konsistentes, wartbares Produktions-Ökosystem.
  
  \item \textbf{Doore et al. (2023)} \cite{Doore2023MDS}.\\
  \emph{Beschreibung:} Evaluationsstudie zu universell zugänglichen Darstellungen von Diagrammen; vergleicht visuelle, taktile und auditive Modalitäten in der Lehre.\\
  \emph{Begründung:} Untermauert den multimodalen Ansatz (Diagramme, Grafiken, Daten) und liefert didaktische Begründung für Multimodalität im CS/EE-Kontext.
  
  \item \textbf{Sharif et al. (2021)} \cite{Sharif2021DataVis}.\\
  \emph{Beschreibung:} Zeigt deutliche Defizite bei Screen-Reader-Zugriff auf Web-Datenvisualisierungen (fehlende Detailtiefe, mangelnde Interaktion).\\
  \emph{Begründung:} Liefert die Problemfolie für das gesamte Paper: Warum zugängliche Diagramme/Grafiken/Daten notwendig sind und welche Fehler zu vermeiden sind.
\end{enumerate}

\section{Related Work 1: Transformable Core Content}
\begin{enumerate}
  \item \textbf{Austin et al. (2024)} \cite{Austin2024PreTeXt}.\\
  \emph{Beschreibung:} XML-basiertes \emph{Single-Source}-Authoring, das in viele barrierefreie Zielformate transformiert; inkl.\ Braille und taktiler Grafiken.\\
  \emph{Begründung:} Fundament für euren \emph{Transformable Core Content}; belegt konsistente Content$\rightarrow$Format-Pipelines in der Praxis.
  
  \item \textbf{Seo \& McCurry (2019)} \cite{Seo2019RMarkdown}.\\
  \emph{Beschreibung:} Praxisbericht zu R~Markdown/Pandoc als zugängliche Alternative; generiert aus einer Quelle u.\,a.\ HTML/PDF und erleichtert Screen-Reader-Navigation.\\
  \emph{Begründung:} Stützt das Markup-first-Narrativ und liefert Referenz für textbasierte Autoren-Workflows mit robusten Exports.
  
  \item \textbf{Armano et al. (2018)} \cite{Armano2018LaTeXPDF}.\\
  \emph{Beschreibung:} Toolchain für zugängliche LaTeX-PDFs mit korrekt getaggten Formeln.\\
  \emph{Begründung:} Ergänzt die Pipeline um das Mathe-Thema (wichtig für CS/EE); zeigt, wie aus Textquellen PDF/UA-nahe Ausgaben entstehen können.
\end{enumerate}

\section{Related Work 2: Accessibility of Diagrams}
\begin{enumerate}
  \item \textbf{Carruthers et al. (2023)} \cite{Carruthers2023PlantUML}.\\
  \emph{Beschreibung:} Erfahrungsbericht zu PlantUML im CS-Kurs; blind studierende Person konnte Diagramme per Text erstellen/bearbeiten und im Team arbeiten.\\
  \emph{Begründung:} Belegt \emph{textuelle Modellierung} als Zugänglichkeits-Pfad; direkt anschlussfähig an eure UML-/PlantUML-Workflows.
  
  \item \textbf{Petrausch et al. (2016)} \cite{Petrausch2016UML}.\\
  \emph{Beschreibung:} Gestaltungsleitlinien für \emph{textuelle} UML-Notationen (Syntaxkürze, Gruppierung, Symbolik).\\
  \emph{Begründung:} Liefert methodische Basis und konkrete Designregeln für eure Diagramm-Beschreibungen.
  
  \item \textbf{Serin \& Romeo (2022)} \cite{Serin2022Diagrams}.\\
  \emph{Beschreibung:} Nutzerstudie mit taktil/akustischer Exploration von UML-Klassendiagrammen (Touchscreen + VoiceOver).\\
  \emph{Begründung:} Stützt die Kombination aus textueller Modellierung und taktiler/akustischer Ergänzung zur Vermittlung räumlicher Struktur.
\end{enumerate}

\section{Related Work 3: Accessibility of Graphics}
\begin{enumerate}
  \item \textbf{Wabi{\'n}ski et al. (2022)} \cite{Wabinski2022TactileMaps}.\\
  \emph{Beschreibung:} Konsolidierte Richtlinien für taktile Darstellungen (Symbole, Linien, Abstände, Beschriftungen).\\
  \emph{Begründung:} Universell auf STEM-Grafiken übertragbar (Legibilität \& Vereinfachung) und damit direkt nutzbar für eure Design-Checks.
  
  \item \textbf{Barros et al. (2023)} \cite{Barros2023Tactile3D}.\\
  \emph{Beschreibung:} Empirische Parameter für 3D-gedruckte, fühlbare Grafiken (Linienbreite/-höhe, Texturen).\\
  \emph{Begründung:} Liefert handfeste Spezifikationen für robuste, gut unterscheidbare Taktilelemente in euren Workflows.
  
  \item \textbf{Clepper et al. (2025)} \cite{Clepper2025Tactile}.\\
  \emph{Beschreibung:} Interviews mit BLV-Taktikerstellenden; Bedarf an alltagstauglichen Tools und umfassender Abbildung verschiedenster Grafiken.\\
  \emph{Begründung:} Verankert Qualitätskriterien (Vollständigkeit, Genauigkeit) aus Nutzerperspektive—passend für euren Bildungs-Kontext.
\end{enumerate}

\section{Related Work 4: Accessibility of Numerical Data}
\begin{enumerate}
  \item \textbf{Sharif et al. (2021)} \cite{Sharif2021DataVis}.\\
  \emph{Beschreibung:} Screen-Reader-Erfahrungen mit Web-Charts; dokumentiert fehlende Detailtiefe und Interaktion.\\
  \emph{Begründung:} Motiviert Tabellen-Alternativen, Pattern/Shape-Mapping und interaktive Zugriffe in eurem Pattern.
  
  \item \textbf{Holloway et al. (2022)} \cite{Holloway2022InfoSonics}.\\
  \emph{Beschreibung:} \emph{Infosonics} koppelt Sonifikation und Sprache; bessere Trend- und Ausreißer-Erkennung für BLV-User.\\
  \emph{Begründung:} Belegt Nutzen multimodaler Audiozugänge—direkt übertragbar auf eure Lehrbeispiele.
  
  \item \textbf{Alam et al. (2023)} \cite{Alam2023SeeChart}.\\
  \emph{Beschreibung:} Sprach-/NLQ-Interface zur explorativen Abfrage von Diagrammdaten.\\
  \emph{Begründung:} Stützt interaktive, text-/sprachbasierte Exploration (Drill-down) als Ergänzung zu statischen Beschreibungen.
\end{enumerate}

\section{Related Work 5: Graphics with Generated Explanation}
\begin{enumerate}
  \item \textbf{Leotta et al. (2023)} \cite{Leotta2023AutoCaptioning}.\\
  \emph{Beschreibung:} Evaluationsstudie automat.\ Bildbeschreibungssysteme für Web-Zugänglichkeit; Qualitätsgrenzen \& Redaktionsbedarf.\\
  \emph{Begründung:} Exakte Evidenz für den \emph{Human-in-the-Loop}-Schritt in eurem Pattern (Qualitätssicherung von generierten Beschreibungen).
  
  \item \textbf{Edwards et al. (2023)} \cite{Edwards2023AltTextProcess}.\\
  \emph{Beschreibung:} Rollen \& Prozesse bei Alt-Texterstellung (Autoren, Redakteur:innen, QA); Praxisleitlinien.\\
  \emph{Begründung:} Blaupause für eure redaktionellen Workflows (Zuständigkeiten, Review, Freigabe).
  
  \item \textbf{Jung et al. (2021)} \cite{Jung2021AltText}.\\
  \emph{Beschreibung:} Anforderungen an Beschreibungen von Datenvisualisierungen aus BLV-Sicht; welche Inhalte sind nötig.\\
  \emph{Begründung:} Liefert konkrete Qualitätskriterien (Genauigkeit, Spezifität, Bezug zu Daten) für eure generierten/editierten Erklärtexte.
\end{enumerate}

\section*{Literatur}
\bibliographystyle{plain}
\bibliography{refs}

\end{document}
